The additional objective comparisons use opened data and do not constitute a new independent test. The choice of learning objective does not reveal a stable extra-window gain in the source cohort. The matched four-cell retrospective analysis contains 54,837 cases in 144 components. Added windows change retention by -0.68 to 0.07 percentage points with quantile learning and 0.02 to 0.12 with direct failure-probability learning. The learning objective does not recover a stable five-point representation advantage. Earlier-local-input and full-procedure results, including their uncertainty, are retained in SI; no non-significant result is declared proof of equivalence.

Across 1,000 full-procedure refits, retrospective window increments average 0.35 points [-4.80, 5.43] for quantile learning and 0.70 [-4.70, 5.82] for failure-probability learning. The wide intervals permit both harm and practically relevant benefit; superiority and equivalence remain unestablished.
